\diarytitle{0028}{ラグランジュ方程式の媒介変数表示による一般解と任意定数の決定（y = 2xy' + (y')^2）}

ラグランジュ方程式 $y = xg(y') + f(y')$（$g(p) \not\equiv p$）は $p = y'$ とおき、$x$ を $p$ の関数とみなして両辺を $p$ で微分することで、$x$ に関する1階線形微分方程式に帰着させて解く。

$p = y'$ とおくと $y = xg(p) + f(p)$ の形（ラグランジュ方程式）で、$g(p) = 2p$, $f(p) = p^{2}$、$g'(p) = 2$, $f'(p) = 2p$ である。$p - g(p) = p - 2p = -p$ に注意して
\[
(p - g(p))\,dx = \bigl[x g'(p) + f'(p)\bigr]\,dp
\quad\Longrightarrow\quad
-p\,dx = (2x + 2p)\,dp
\]
\[
\frac{dx}{dp} = -\frac{2x + 2p}{p} = -\frac{2}{p}x - 2
\quad\Longrightarrow\quad
\frac{dx}{dp} + \frac{2}{p}\,x = -2
\]
これは $x$ に関する1階線形微分方程式である。積分因子は $\mu(p) = \exp\left(\displaystyle\int \frac{2}{p}\,dp\right) = p^{2}$ なので
\[
\frac{d}{dp}\bigl(x\,p^{2}\bigr) = -2p^{2}
\quad\Longrightarrow\quad
x\,p^{2} = -\frac{2}{3}p^{3} + C
\]
\[
\boxed{\; x(p) = -\frac{2}{3}p + \frac{C}{p^{2}} \;}
\qquad (p \neq 0)
\]
これを $y = 2xp + p^{2}$ に代入すると
\[
y = 2p\left(-\frac{2}{3}p + \frac{C}{p^{2}}\right) + p^{2}
= -\frac{4}{3}p^{2} + \frac{2C}{p} + p^{2}
\]
\[
\boxed{\; y(p) = -\frac{1}{3}p^{2} + \frac{2C}{p} \;}
\]
（検算: $\dfrac{dx}{dp} = -\dfrac{2}{3} - \dfrac{2C}{p^{3}}$, $\dfrac{dy}{dp} = -\dfrac{2}{3}p - \dfrac{2C}{p^{2}} = p\left(-\dfrac{2}{3} - \dfrac{2C}{p^{3}}\right) = p\,\dfrac{dx}{dp}$ より $\dfrac{dy}{dx} = p$ を満たす。また $2xp + p^{2} = 2p\left(-\dfrac23 p + \dfrac{C}{p^2}\right)+p^2 = -\dfrac43p^2+\dfrac{2C}{p}+p^2 = -\dfrac13 p^2+\dfrac{2C}{p} = y$ となり、原方程式を満たす。）

条件: $p = 1$ のとき $x = -\dfrac{1}{3}$, $y = \dfrac{1}{3}$ を通る。$x(1) = -\dfrac{2}{3} + C = -\dfrac{1}{3}$ より
\[
\boxed{\; C = \frac{1}{3} \;}
\]
（検算: $y(1) = -\dfrac{1}{3} + 2 \cdot \dfrac{1}{3} = -\dfrac{1}{3} + \dfrac{2}{3} = \dfrac{1}{3}$ となり、条件を満たす。）
